\section{Introduction}
\label{sec:introduction}

% Hook: The problem
Large language model (LLM) agents face a fundamental constraint: the
context window is finite, but experience is not. An agent that
converses across hundreds of sessions accumulates knowledge,
preferences, and relationships that cannot fit in a single prompt.
Memory---the ability to store, retrieve, and manage this accumulation
over time---is now a first-class architectural concern.

% Scale of the problem
The urgency is measurable. OpenClaw, the most widely deployed AI coding
agent, injects approximately 35,600 tokens of workspace files into
every message, 93.5\% of which is irrelevant to the current
query~\cite{openclaw9157}. Heavy users report costs of \$3,600 per
month in token waste alone~\cite{zilliz2026b}. Community tools that
replace full-context injection with ranked retrieval report 70--96\%
token savings~\cite{levine2026}, demonstrating that the problem is not
the absence of memory but the absence of \emph{selective} memory.

% Explosion of systems
Over 90 memory systems have emerged since Park~\etal's seminal
Generative Agents work in 2023~\cite{park2023generative}. These span
from flat-file approaches (MEMORY.md) through knowledge-graph engines
(Zep/Graphiti~\cite{dempsey2025graphiti},
HippoRAG~\cite{gutierrez2024hipporag}) to cognitively inspired
architectures (SYNAPSE~\cite{jiang2026synapse},
LightMem~\cite{fang2025lightmem}). The pace of development outstrips
synthesis: six survey papers appeared in 2024--2026
alone~\cite{sumers2024coala, gao2023rag, zhang2024survey, hu2025memory,
wu2025memory, du2025memory}, yet none covers the full landscape from
production deployments through practitioner-invented workarounds.

% Gap this paper fills
This paper provides the first comprehensive taxonomy that bridges three
communities: (1) \textbf{academic systems} published at NeurIPS, ICML,
ICLR, ACL, and AAAI; (2) \textbf{production systems} deployed by
companies (Mem0, Zep, Supermemory, Letta); and (3)
\textbf{practitioner workarounds} invented by users of file-based
memory systems (decision.md, working-context.md, 12-layer
architectures). We organize these along four axes:

\begin{enumerate}[leftmargin=*, nosep]
    \item \textbf{Structure:} What memory types does the system
    support? Mapped to the CoALA taxonomy~\cite{sumers2024coala}.
    \item \textbf{Retrieval:} How does the system find relevant
    memories? Mapped to the RAG taxonomy~\cite{gao2023rag}.
    \item \textbf{Lifecycle:} How does the system manage memory over
    time? Formation, consolidation, forgetting, transformation.
    \item \textbf{Preference emergence:} How do behavioral patterns
    crystallize from accumulated interaction?
\end{enumerate}

% The interdisciplinary thesis
We ground this taxonomy in two traditions that have studied memory for
far longer than computer science. Cognitive neuroscience provides
models of Hebbian learning (LTP/LTD)~\cite{hebb1949, bliss1973},
complementary learning systems~\cite{mcclelland1995}, spreading
activation~\cite{collins1975}, and dual-strength
forgetting~\cite{bjork1992}. Yogac\=ara Buddhist psychology, developed
over 1,500 years, provides a theory of \emph{storehouse consciousness}
(\emph{alaya-vij\~n\=ana}) with \emph{seeds} (\emph{b\=ija}) that
strengthen through \emph{perfuming} (\emph{v\=asan\=a}) and transform
through \emph{\=a\'sraya-par\=av\d{r}tti}. We show that both
traditions converge on the same insight: \textbf{memory is a dynamic
process that reshapes through use, not a static database that grows
through append}.

% Contributions
Our contributions are:
\begin{enumerate}[leftmargin=*, nosep]
    \item A four-axis taxonomy covering 90+ memory systems across
    academic, production, and practitioner categories.
    \item Comparative matrices mapping systems to CoALA memory
    modules, RAG retrieval categories, and lifecycle operations.
    \item Documentation of the MEMORY.md problem and five convergent
    design principles the community has independently discovered.
    \item A Yogac\=ara framework that maps Buddhist psychology concepts
    to concrete memory architecture components.
    \item An open research agenda identifying underserved areas:
    principled forgetting, preference emergence, adaptive retrieval,
    and cross-memory consolidation.
\end{enumerate}
